\chapter{Poincar\'e 群、Casimir 与 Wigner 分类}
\label{chap:lorentz-poincare}

$(j_L,j_R)$ 说明一个局域场的分量在 Lorentz 变换下怎样混合，却不能直接充当粒子的自旋标签：例如四矢量场限制到旋转子群时含有多个角动量分量，而有质量自旋一粒子的物理态只保留三个极化。要分类单粒子态，必须把平移也纳入对称群，并要求它在 Hilbert 空间上幺正实现。

平移生成元的联合谱先把态限制到一个动量轨道；保持某个标准动量不变的小群再作用于内部指标。这样，质量来自轨道，内禀自旋或螺旋度来自小群表示。Pauli--Lubanski 矢量把这两层写成与参考系无关的 Casimir 算符。

\section{Poincar\'e 群及其 Casimir 算符}

Poincar\'e 变换由平移 $a\in\mathbb R^{3,1}$ 与 $\Lambda\in SO^+(3,1)$ 组成，对时空作用为 $x\mapsto\Lambda x+a$。连续群是半直积 $\mathcal P_+^\uparrow=\mathbb R^{3,1}\rtimes SO^+(3,1)$，群乘法为
\[
(a,\Lambda)(b,\Gamma)=(a+\Lambda b,\Lambda\Gamma).
\]
它不是直积，因为 Lorentz 变换会改变平移四矢量。Poincar\'e 代数式 \eqnrefs{eq:poincare-algebra-1} 与式 \eqnrefs{eq:poincare-algebra-2} 正是这一半直积结构的 Lie 代数表述。

令 $P_\mu$ 为平移生成元，则 Lie 代数除 Lorentz 部分外还满足
\begin{align}
[P_\mu,P_\nu]&=0,
\label{eq:poincare-algebra-1}\\
[M_{\mu\nu},P_\rho]
&=\ii(\eta_{\nu\rho}P_\mu-\eta_{\mu\rho}P_\nu).
\label{eq:poincare-algebra-2}
\end{align}
第一 Casimir 是 $P^2=P_\mu P^\mu$。第二个 Casimir 由 Pauli--Lubanski 四矢量（式 \eqnrefs{eq:pauli-lubanski}，取 $\epsilon^{0123}=+1$ 的 Levi--Civita 约定）
\begin{equation}
W^\mu\equiv\frac12\epsilon^{\mu\nu\rho\sigma}P_\nu M_{\rho\sigma}
\label{eq:pauli-lubanski}
\end{equation}
构成。它满足 $W\cdot P=0$，而 $W^2$ 与所有 Poincar\'e 生成元对易。对质量 $m>0$ 的不可约表示，可取静止动量 $k=(m,\vb{0})$；此时 $W^0=0$、$\vb{W}=m\vb{J}$，所以
\begin{equation}
W^2=-m^2j(j+1).
\label{eq:W2-massive}
\end{equation}
式 \eqnrefs{eq:W2-massive} 说明 $P^2=m^2$ 与 $W^2$ 分别读出质量和自旋。

\section{单粒子态与 Wigner 分类}
\label{sec:wigner-classification}

先选一个标准动量 $k$，再考察保持它不变的 Lorentz 子群 $G_k=\{\Lambda\in SO^+(3,1)\mid\Lambda k=k\}$，称为 $k$ 的小群。动量轨道上的任意 $p$ 可写为 $p=L(p)k$。于是 Poincar\'e 表示由小群表示诱导出来：群先改变动量，再由
\[
R(\Lambda,p)=L(\Lambda p)^{-1}\Lambda L(p)\in G_k
\]
在内部指标上作用。

对正能有质量轨道 $p^2=m^2$、$p^0>0$，可取 $k=(m,\vb{0})$，小群为 $SO(3)$；若同时允许半整数自旋，应使用其覆盖群 $SU(2)$。表示空间可写成
\[
\mathcal H_{m,j}=L^2(\mathcal O_m,\dd\mu_m)\otimes\mathbb C^{2j+1},
\qquad
\dd\mu_m(p)=\frac{\dd^3p}{2E_{\vb{p}}},
\]
而群作用为
\[
(U(a,\Lambda)\psi)_\sigma(p)
=\ee^{\ii a\cdot p}
\sum_{\sigma'}D^{(j)}_{\sigma\sigma'}\!\left(R(\Lambda,\Lambda^{-1}p)\right)
\psi_{\sigma'}(\Lambda^{-1}p).
\]
不变测度保证 $U$ 幺正，小群不可约性保证诱导表示不可约。

对正能无质量轨道，可取 $k=(\kappa,0,0,\kappa)$，小群同构于 $ISO(2)$。若其二维平移子群作用平凡，表示由单个螺旋度 $h$ 标记，并满足 $W^\mu=hP^\mu$；宇称若是对称性，通常把 $h$ 与 $-h$ 配成一对。若平移子群非平凡，则得到连续自旋表示，它含无穷多个螺旋度，不能被误写成普通有限螺旋度粒子。

空间反演在 Dirac 场上的常用实现是
\[
\mathsf P\,\Psi(t,\vb{x})\,\mathsf P^{-1}
=\eta_P\gamma^0\Psi(t,-\vb{x}),
\]
其中 $|\eta_P|=1$ 是内禀宇称相位。时间反演必须反幺正；在常用表象中可取
\[
\mathsf T\,\Psi(t,\vb{x})\,\mathsf T^{-1}
=\eta_T\,\ii\gamma^1\gamma^3\Psi(-t,\vb{x}),
\]
并同时复共轭数系数。相位约定可以改变，反幺正性及其对角动量的反号作用不能省略。

\begin{derivation}[诱导表示中的 Wigner 转动]
定义标准态 $|p,\sigma\rangle=U(L(p))|k,\sigma\rangle$。对 Lorentz 变换 $\Lambda$，插入恒等分解
\[
\Lambda L(p)=L(\Lambda p)\bigl[L(\Lambda p)^{-1}\Lambda L(p)\bigr]
=L(\Lambda p)R(\Lambda,p).
\]
方括号保持 $k$ 不变，因而只能通过小群矩阵 $D(R)$ 混合内部指标。平移再给出相位 $\ee^{\ii a\cdot p}$，就得到正文的诱导表示公式。整个分类所需的信息因此是动量轨道和小群不可约表示，而不是 Lorentz 群正则表示的 Plancherel 分解。
\end{derivation}

有限维 Lorentz 表示 $(j_L,j_R)$ 描述局域场的分量结构，通常并不幺正；Poincar\'e 群的不可约幺正表示描述单粒子 Hilbert 空间。有质量轨道的小群是 $SU(2)$，给出自旋 $j$；无质量轨道的小群是 $ISO(2)$，其平移部分平凡时给出螺旋度。Clifford 代数把左右 Weyl 表示组合成 Dirac 旋量，而宇称交换两种手征。这些结论共同把时空几何、场的变换律和粒子量子数连成一条严格的表示论链。
